\documentclass[10pt,a4paper]{amsart}
\usepackage[margin=19mm]{geometry}
\usepackage{amsmath,amssymb,mathtools}
\usepackage[scr=rsfs,cal=euler]{mathalfa}
\usepackage{xfrac}
\usepackage[unicode,hidelinks,bookmarks=false]{hyperref}
\newcommand{\ie}{i.e.\ }
\newcommand{\cf}{cf.\ }
\newcommand{\resp}{respectively\ }
\newcommand{\Subsection}{\subsection}
\providecommand{\nobreakdash}{\nobreak\hbox{-}}
\setlength{\emergencystretch}{2em}
\multlinegap0pt
\usepackage{bm}
\usepackage{bbm}
\usepackage{dsfont}
\usepackage{xspace}
\usepackage{cancel}
\usepackage{mathtools}
\usepackage{amsthm}
\makeatletter
\@ifundefined{theorem}{\newtheorem{theorem}{Theorem}}{}
\@ifundefined{lemma}{\newtheorem{lemma}[theorem]{Lemma}}{}
\makeatother
\def \Mvec{\mathbf{M}}
\def\Hvec{\mathbf{H}}
\def\Qvec{\mathbf{Q}}
\def\d{\mathrm{d}}
\title[A study of ferronematic thin films including a stray field e]{A study of ferronematic thin films including a stray field energy}
\author{Shilpa Dutta, James Dalby, Apala Majumdar, Anja Schl\"omerkemper}
\date{Source: arXiv:2509.10442v1 (CC BY 4.0)}
\begin{document}
\maketitle

\noindent\emph{Source and licence.} A study of ferronematic thin films including a stray field energy. Version arXiv:2509.10442v1 (CC BY 4.0), distributed under the Creative Commons Attribution 4.0 International licence, as recorded in the arXiv licence metadata for this exact version and shown on the abstract page https://arxiv.org/abs/2509.10442v1. Full rights notice: licenses/arxiv-2509.10442-CC-BY-4.0.txt.\par\smallskip

\noindent\textit{Editorial scope.} Counted unit: the Theorem whose source label is \texttt{thm:2} in the article (source0.tex lines 449--451, source label \texttt{thm:2}), delivered together with the article's own complete original proof of it (source0.tex lines 452--532, one \texttt{proof} environment of 81 lines). No mathematics is written, repaired or re-derived: statement and proof are the article's own text line for line. Required context retained uncounted: en:2 (lines 97--103); def-f (lines 234--236); EL:1 (lines 353--364); EL:5 (lines 353--364); lem1 (lines 369--379). Every definition, display and label of those lines is retained unchanged. Omitted from this package and not redistributed: 1--96, 104--233, 237--352, 365--368, 380--448, 533--1000. Nothing inside a retained excerpt is omitted or altered: every retained body is delivered exactly as the article's lines are written, so no omission receipt of any kind accompanies this package. Citations: the retained excerpts carry no citation command at all, so no bibliography entry of the article is transcribed and no reference is redistributed. Reference adaptation: none was needed and none was made; every reference inside the delivered unit resolves inside it, so no unresolved reference or citation is produced. Printed slips kept literal: no printed word, symbol, hypothesis or number of the counted statement or of its proof was altered. Layout and class: the article's own class is \texttt{article} with packages including geometry, amsthm, titling, mathtools, xspace, mathrsfs, esint, algorithm2e, algorithmic, graphicx, caption, subcaption, bm, amsmath; the wrapper is the lane's pinned \texttt{amsart} editorial wrapper at 10pt on A4, which restates the theorem-like environments the retained text needs and adds the article's own macro definitions that the retained text needs (\texttt{\textbackslash Hvec}, \texttt{\textbackslash Mvec}, \texttt{\textbackslash Qvec}, \texttt{\textbackslash d}), with extra wrapper packages (bm, bbm, dsfont, xspace, cancel, mathtools). The delivered numbering is the unit's own. Assets: no retained span contains a figure environment, a graphics-inclusion command, a PDF-page inclusion, a table or a TikZ picture; no image file is copied into this package, the package declares no asset dependency and no image file is redistributed with it. Licence: version arXiv:2509.10442v1 of this article is distributed under the Creative Commons Attribution 4.0 International licence, as recorded in the arXiv licence metadata for this exact version and shown on the abstract page https://arxiv.org/abs/2509.10442v1. A full rights notice is delivered at licenses/arxiv-2509.10442-CC-BY-4.0.txt. Characters: every retained excerpt is pure ASCII, so the delivered engine, XeLaTeX with its default Latin Modern text font, reports no missing character.
\begin{eqnarray}\label{en:2}
\mathcal{E}\left(\Qvec, \Mvec\right)&=&\int_{\Omega}\frac{l_1}{2}|\nabla \Qvec|^2 + \frac{\xi l_2}{2}|\nabla \Mvec|^2 - \frac{c_1}{2}\Qvec\Mvec\cdot\Mvec 
+ \frac{c_3}{2}\xi\left(M_3\right)^2
\nonumber\\
&& \;\;- \frac{c_2}{2}\Qvec\Hvec_{ext}\cdot\Hvec_{ext} 
- c_3\xi\Mvec\cdot\Hvec_{ext} + \frac{1}{4}\left(\frac{1}{2}|\Qvec|^2 - 1\right)^2 + \frac{\xi}{4}\left(|\Mvec|^2 - 1\right)^2 \,\d x,
\end{eqnarray}

\begin{eqnarray} \label{def-f}
f\left(\Qvec, \Mvec\right)&=&  \frac{1}{4}\left(\frac{1}{2}|\Qvec|^2 - 1\right)^2 + \frac{\xi}{4}\left(|\Mvec|^2 - 1\right)^2 -  \frac{c_1}{2}\Qvec\Mvec\cdot\Mvec.
\end{eqnarray}

\begin{align}
\label{EL:1}
2l_1\Delta Q_{11} &=  Q_{11}\left(Q_{11}^2 + Q_{12}^2 -1\right) - \frac{c_1}{2}\left(M_1^2 - M_2^2\right) - \frac{c_2}{2}\left(H_1^2 - H_2^2\right),\\
\label{EL:2}
2l_1\Delta Q_{12} &=  Q_{12}\left(Q_{11}^2 + Q_{12}^2 - 1\right) - c_1M_1M_2 - c_2H_1H_2,\\
\label{EL:3}
l_2\xi\Delta M_1 &= \xi M_1\left(M_1^2 + M_2^2 + M_3^2 - 1\right) - c_1\left(Q_{11}M_1 + Q_{12}M_2\right) - c_3\xi H_1,\\
\label{EL:4}
l_2\xi\Delta M_2 &= \xi M_2\left(M_1^2 + M_2^2 + M_3^2 - 1\right) - c_1\left(Q_{12}M_1 - Q_{11}M_2\right) - c_3\xi H_2,\\
\label{EL:5}
l_2\xi\Delta M_3 &= \xi M_3\left(M_1^2 + M_2^2 + M_3^2 - 1\right) - c_3\xi \left(H_3-M_3\right).
\end{align}

\begin{lemma}\label{lem1}
Let $\left(\Qvec, \Mvec\right)\in\mathcal{A}\times\mathcal{S}$ be an $H^1$ weak solution of the system of PDEs \eqref{EL:1}--\eqref{EL:5}, where $\Hvec_{ext}\in C\left(\overline{\Omega}; \mathbb R^3\right)$ with $\sup_{\bar{\Omega}}|\Hvec_{ext}|\le k_1$. Then the weak solution $\left(\Qvec, \Mvec\right)$ has an interior local bound in terms of the positive parameters, $c_1, c_2, c_3,$ $k_1$ and $\xi$, i.e., for any compact set $K\subset\Omega$, there exists $\sigma = \sigma\left(c_1, c_2, c_3, \xi, k_1\right)$ such that 
\begin{eqnarray} 
\begin{aligned}
|Q_{11}(z)| &\leq \sigma~~~~~~~~~~~~~~~~~~~~~\forall z\in K,\label{eq:Qmax}\\
|Q_{12}(z)| &\leq \sigma~~~~~~~~~~~~~~~~~~~~~\forall z\in K,\\
|\Mvec(z)|^2 &\leq A c_1 \frac{\sigma}{\xi} + Bc_3k_1 + 1~~\forall z\in K, \label{eq:Mmax}
\end{aligned}
\end{eqnarray}
for some positive constants $A$ and $B$ independent of $c_1, c_2, c_3, \xi, l_1, l_2, k_1$.
\end{lemma}

\begin{theorem}[Uniqueness of minimizer]\label{thm:2}
Let the parameters $c_1,c_2,c_3, \xi>0$ be fixed and $\Hvec_{ext}\in C\left(\overline{\Omega}; \mathbb R^3\right)$  with $\sup_{\bar{\Omega}}|\Hvec_{ext}|\le k_1$. For sufficiently large $l_1$ and $l_2$, there exists a unique minimizer of the ferronematic energy functional $\mathcal{E}$, defined in \eqref{en:2}, in the admissible space $\mathcal{A}\times\mathcal{S}$. 
\end{theorem}

\begin{proof}
On the contrary, we assume that $\left(\Qvec, \Mvec\right)$ and $\left(\overline{\Qvec}, \overline{\Mvec}\right)$ $\in$ $\mathcal{A}\times\mathcal{S}$ are two distinct minimizers of $\mathcal{E}$. By an elementary calculation and the definition of $f$ in \eqref{def-f}, we obtain
\begin{eqnarray}
& &\mathcal{E}\left(\frac{\Qvec+\overline{\Qvec}}{2}, \frac{\Mvec+\overline{\Mvec}}{2}\right)\nonumber\\
&=&\int_{\Omega}\frac{l_1}{2}\bigg|\nabla\left(\frac{\Qvec+\overline{\Qvec}}{2}\right)\bigg|^2 + \frac{\xi l_2}{2}\bigg|\nabla\left(\frac{\Mvec+\overline{\Mvec}}{2}\right)\bigg|^2 + f\left(\frac{\Qvec+\overline{\Qvec}}{2}, \frac{\Mvec+\overline{\Mvec}}{2}\right)+\frac{c_3}{2}\xi\left(\frac{M_3+\overline{M_3}}{2}\right)^2\nonumber\\
& &- \frac{c_2}{2}\left(\frac{\Qvec+\overline{\Qvec}}{2}\right)\Hvec_{ext}\cdot\Hvec_{ext} 
- c_3\xi\left(\frac{\Mvec+\overline{\Mvec}}{2}\right)\cdot\Hvec_{ext}\;\d x\nonumber\\
&=&\frac{1}{2}\mathcal{E}\left(\Qvec, \Mvec\right) + \frac{1}{2}\mathcal{E}\left(\overline{\Qvec}, \overline{\Mvec}\right) + \int_{\Omega}f\left(\frac{\Qvec+\overline{\Qvec}}{2}, \frac{\Mvec+\overline{\Mvec}}{2}\right)-\frac{1}{2}\left(f\left(\Qvec, \Mvec\right) + f\left(\overline{\Qvec}, \overline{\Mvec}\right)\right)\nonumber\\
& &- \frac{l_1}{8}\left(\nabla\Qvec-\nabla\overline{\Qvec}\right)^2-\frac{l_2\xi}{8}\left(\nabla\Mvec-\nabla\overline{\Mvec}\right)^2 - \frac{c_3\xi}{8}\left(M_3-\overline{M_3}\right)^2\d x\nonumber\\
&\le&\frac{1}{2}\left(\mathcal{E}\left(\Qvec, \Mvec\right) + \mathcal{E}\left(\overline{\Qvec}, \overline{\Mvec}\right)\right)+\int_{\Omega}f\left(\frac{\Qvec+\overline{\Qvec}}{2}, \frac{\Mvec+\overline{\Mvec}}{2}\right)-\frac{1}{2}\left(f\left(\Qvec, \Mvec\right) + f\left(\overline{\Qvec}, \overline{\Mvec}\right)\right)\d x\nonumber\\
& &- \frac{l_1}{8c_p}||\Qvec-\overline{\Qvec}||^2_{L^2\left(\Omega; \mathbb R^{3\times 3}\right)} -\frac{l_2\xi}{8c_p}||\Mvec-\overline{\Mvec}||^2_{L^2\left(\Omega; \mathbb R^3\right)} - \frac{c_3\xi}{8}||M_3-\bar{M_3}||^2_{L^2\left(\Omega; \mathbb R\right)},\nonumber\\ \label{eq:un1}
\end{eqnarray}
where we have used the Poincar\'e inequality 
with constant $c_p$ (independent of $\Omega$) in the last line. Since $\Qvec-\overline{\Qvec}, \Mvec-\overline{\Mvec}\in \overline{C^{\infty}_{c}\left(\Omega\right)}$, there exists a compact set $K\subset\Omega$ such that
\begin{eqnarray}
\begin{aligned}
\Qvec-\overline{\Qvec}=0 \mbox{ a.e. in } \Omega\setminus K,\label{eq:tay}\\
\Mvec-\overline{\Mvec}=0 \mbox{ a.e. in } \Omega\setminus K.
\end{aligned}
\end{eqnarray}
Moreover, $f\left(\Qvec, \Mvec\right)$ is $C^{\infty}$ w.r.t $Q_{ij}, M_j$ and we deal with the following two scenarios.\newline\\
\textbf{Case I.} Assume that $|\left(\Qvec, \Mvec\right)-\left(\overline{\Qvec}, \overline{\Mvec}\right)|\geq 1$. We know that $f$ is locally Lipschitz. As a consequence, we have
\begin{eqnarray}
& &f\left(\frac{\Qvec+\overline{\Qvec}}{2}, \frac{\Mvec+\overline{\Mvec}}{2}\right)-\frac{1}{2}\left(f\left(\Qvec, \Mvec\right) + f\left(\overline{\Qvec}, \overline{\Mvec}\right)\right)\nonumber\\
&\le&\lambda|\left(\Qvec, \Mvec\right)-\left(\overline{\Qvec}, \overline{\Mvec}\right)|^2,\label{eq:case1}
\end{eqnarray}
where $\lambda$ is the Lipschitz constant and $\lambda=\max_{\left(\Qvec, \Mvec\right)}|f'\left(\Qvec, \Mvec\right)|$. Applying Lemma~\ref{lem1}, the first partial derivatives of $f$ over the compact set $K$ can be estimated as 
\begin{eqnarray}
\Big|\dfrac{\partial f}{\partial Q_{1j}}\Big|&\le&|Q_{ij}|\left|\tfrac12|\Qvec|^2-1\right|+c_1|\Mvec|^2,\nonumber\\
&\le& \sigma^3 + \sigma + c_1\left(A c_1 \frac{\sigma}{\xi} + Bc_3k_1 + 1\right):=a_1,\nonumber\\
\label{eq:first_de1}\\
\Big|\dfrac{\partial f}{\partial M_{i}}\Big|&\le&\xi|M_i|\left||\Mvec|^2-1\right|+2c_1\left(\tfrac12|\Qvec|^2+|\Mvec|^2\right)\nonumber\\
&\le&\xi\sqrt{\left(A c_1 \frac{\sigma}{\xi} + Bc_3k_1 + 1\right)}\left(A c_1 \frac{\sigma}{\xi} + Bc_3k_1 + 2\right)+2c_1\left(\sigma^2+ A c_1 \frac{\sigma}{\xi} + Bc_3k_1  + 1\right)\nonumber\\
&:=&a_2.\nonumber\\
\label{eq:first_de2}
\end{eqnarray}
\textbf{Case II.} Assume that $|\left(\Qvec, \Mvec\right)-\left(\overline{\Qvec}, \overline{\Mvec}\right)| < 1$.\newline
By Taylor's expansion followed by an elementary calculation, we obtain that
\begin{eqnarray}\label{eq:un2'}
\lefteqn{
f\left(\frac{\Qvec+\overline{\Qvec}}{2}, \frac{\Mvec+\overline{\Mvec}}{2}\right)-\frac{1}{2}\left(f\left(\Qvec, \Mvec\right) + f\left(\overline{\Qvec}, \overline{\Mvec}\right)\right)}\nonumber\\
&\le& \frac{1}{4}\bigg|f^{''}\left(\frac{\Qvec+\overline{\Qvec}}{2}, \frac{\Mvec+\overline{\Mvec}}{2}\right)\bigg||\left(\Qvec, \Mvec\right)-\left(\overline{\Qvec}, \overline{\Mvec}\right)|^2 \nonumber\\
&&+\frac{1}{16}\bigg|f^{iv}\left(\frac{\Qvec+\overline{\Qvec}}{2}, \frac{\Mvec+\overline{\Mvec}}{2}\right)\bigg||\left(\Qvec, \Mvec\right)-\left(\overline{\Qvec}, \overline{\Mvec}\right)|^4\nonumber\\
&\le& \frac{1}{4}\bigg|f^{''}\left(\frac{\Qvec+\overline{\Qvec}}{2}, \frac{\Mvec+\overline{\Mvec}}{2}\right)\bigg||\left(\Qvec, \Mvec\right)-\left(\overline{\Qvec}, \overline{\Mvec}\right)|^2 \nonumber\\ && +\frac{1}{8}\bigg|f^{iv}\left(\frac{\Qvec+\overline{\Qvec}}{2}, \frac{\Mvec+\overline{\Mvec}}{2}\right)\bigg||\left(\Qvec, \Mvec\right)-\left(\overline{\Qvec}, \overline{\Mvec}\right)|^2.\nonumber\\
\end{eqnarray}
Applying Lemma~\ref{lem1}, the estimates for the second partial derivatives of $f\left(\Qvec, \Mvec\right)$ on the compact set $K\subset\Omega$, are obtained as 
\begin{eqnarray}
\begin{aligned}
\bigg|\frac{\partial^2 f}{\partial Q_{1i}\partial Q_{1j}}\bigg|&\le A_1\sigma^2:=a'_1,\\
\biggl|\frac{\partial^2 f}{\partial M_{i}\partial M_{j}}\biggr|&\le A_2c_1\sigma+B_2\xi c_3k_1+C_2\xi:=a'_2,\\
\biggl|\frac{\partial^2 f}{\partial Q_{1i}\partial M_{j}}\biggr|&\le A_3c_1\sqrt{Ac_1\frac{\sigma}{\xi}+Bc_3k_1+1 }:=a'_3,\\
\end{aligned}\label{eq:estimates}
\end{eqnarray}
where $A_1, A_2, A_3, B_2$ and $C_2$ are some positive constants independent of $c_1, c_2, c_3, \xi, k_1, l_1, l_2$.\newline\\
Due to the estimates \eqref{eq:case1}, \eqref{eq:first_de1}, \eqref{eq:first_de2}, \eqref{eq:un2'}, \eqref{eq:estimates}, and using \eqref{eq:tay}, we get
\begin{eqnarray}\label{eq:un2}
& &\int_{\Omega}f\left(\frac{\Qvec+\overline{\Qvec}}{2}, \frac{\Mvec+\overline{\Mvec}}{2}\right)-\frac{1}{2}\left(f\left(\Qvec, \Mvec\right) + f\left(\overline{\Qvec}, \overline{\Mvec}\right)\right)\d x\nonumber\\
&\le& \left(a_1+a'_1+a'_3+\zeta\right)||\Qvec-\overline{\Qvec}||^2_{L^2\left(\Omega; \mathbb R^{3\times 3}\right)} + \left(a_2+a'_2+a'_3+\zeta\right)||\Mvec-\overline{\Mvec}||^2_{L^2\left(\Omega;\mathbb R^3\right)},
\end{eqnarray}
where the constant $\zeta := \bigg|f^{iv}\left(\frac{\Qvec+\overline{\Qvec}}{2}, \frac{\Mvec+\overline{\Mvec}}{2}\right)\bigg|$. 
Now, using \eqref{eq:un2} in \eqref{eq:un1}, we have
\begin{eqnarray}\label{eq:un4}
\lefteqn{\mathcal{E}\left(\frac{\Qvec+\overline{\Qvec}}{2}, \frac{\Mvec+\overline{\Mvec}}{2}\right)}\nonumber\\
&\le&\frac{1}{2}\mathcal{E}\left(\Qvec, \Mvec\right) + \frac{1}{2}\mathcal{E}\left(\bar{\Qvec}, \overline{\Mvec}\right)+\left(a_1+a'_1+a'_3+\zeta-\frac{l_1}{8c_p}\right)||\Qvec-\overline{\Qvec}||^2_{L^2\left(\Omega; \mathbb R^{3\times 3}\right)}\nonumber\\
&&+ \left(a_2+a'_2+a'_3+\zeta-\frac{l_2\xi}{8c_p}\right)||\Mvec-\overline{\Mvec}||^2_{L^2\left(\Omega; \mathbb R^3\right)} - \frac{c_3\xi}{8}||M_3-\overline{M_3}||^2_{L^2\left(\Omega; \mathbb R\right)}.
\end{eqnarray}
From \eqref{eq:un4}, it can be observed that
\begin{eqnarray}\label{eq:un5}
\begin{aligned}
\mathcal{E}\left(\frac{\Qvec+\overline{\Qvec}}{2}, \frac{\Mvec+\overline{\Mvec}}{2}\right)<\frac{1}{2}\mathcal{E}\left(\Qvec, \Mvec\right) + \frac{1}{2}\mathcal{E}\left(\overline{\Qvec}, \overline{\Mvec}\right),
\end{aligned}
\end{eqnarray}
when $l_1, l_2>l'\left(c_1, c_2, c_3, \xi, k_1\right)=\max \Big\{8c_p\left(a_1+a'_1+a'_3+\zeta\right), \frac{8c_p}{\xi}\left(a_2+a'_2+a'_3+\zeta\right)\Big\}$.\newline
To obtain a contradiction, we assume that $\left(\Qvec, \Mvec\right)\ne \left(\overline{\Qvec}, \overline{\Mvec}\right)\in\mathcal{A}\times\mathcal{S}$ are minimizers of $\mathcal{E}$ for $l_1,l_2\in\left(l', \infty\right)$. Then the mapping $\phi:[0,1]\rightarrow\mathbb R$, defined by
\begin{eqnarray}
\begin{aligned}
\phi\left(\lambda\right)=\mathcal{E}\left(\lambda\Qvec+\left(1-\lambda\right)\overline{\Qvec}, \lambda\Mvec + \left(1-\lambda\right)\overline{\Mvec}\right),
\end{aligned}
\end{eqnarray}
is $C^1$ and $\phi'\left(\lambda\right)=0$ at $\lambda=0,1$. This contradicts the strictly convex nature c.f.\ \eqref{eq:un5} of the energy functional $\mathcal{E}$. Hence Theorem~\ref{thm:2} follows.
\end{proof}

\end{document}
